\begingroup
\renewcommand{\arraystretch}{1.45}
\fancyhead[L]{\small LEMNISKATISCHE FUNKTIONEN II.}
\newcommand{\GaussIIITranscriptionDigits}[1]{\text{\ttfamily #1}}
\SourceLeaf{435}

\par\medskip\Needspace{5\baselineskip}\pdfbookmark[0]{Sammlung von Rechnungen}{sammlung-rechnungen}\centerline{[17.]}

\begin{center}\textit{Sammlung von Rechnungen,}\end{center}
hauptsächlich solcher, bei denen meine Methoden zur Bestimmung der Faktoren
großer Zahlen und die \textsc{Wolframschen} Logarithmentafeln benutzt wurden.

\Needspace{26\baselineskip}
\begin{center}\textit{Erste Rechnung für $e^{-\pi}=A$}\end{center}

Aus einer vorläufigen Näherung war bereits bekannt, dass $A$ bis auf die
$14^{\mathrm{te}}$ Ziffer gleich $0,0432139182\ 6377$ ist; folglich war $11.10^{10}A$ sehr genau gleich $4753531008$
$=128.243.152827$. Es stellte sich also die Frage, ob sich die Zahl $152827$ noch weiter in einfache
Faktoren zerlegen lässt.

\Needspace{26\baselineskip}
Die Division durch kleine Primzahlen gelang nicht; man musste daher
auf besondere Verfahren zurückgreifen. Verglich man die Zahl selbst und
verschiedene ihrer Vielfachen mit den nächstgelegenen Quadraten, so fand sich
unter anderem
\[
552^2\equiv950,\quad677^2\equiv152\quad\mathrm{mod.}\ 152827
\]
woraus man sofort schloss, dass
\[
1104^2\equiv3800,\quad3385^2\equiv3800
\]
und somit die Zahl $152827$ keine Primzahl ist, weil sonst die Quadrate zweier
Zahlen, die beide kleiner als die Hälfte dieser Zahl sind, unmöglich kongruent
sein könnten (\textit{Disqu. Arr. art.}\quad.). Mit der in diesem Werk gelehrten Methode
(\textit{art.}\quad) fand man nun die Faktoren der Zahl $152827$, nämlich $67.2281$.

\Needspace{26\baselineskip}
Man war also sicher, dass sich der hyperbolische Logarithmus von
\[
N=128.243.67.2281.11^{-1}.10^{-10}
\]
aus den vorhandenen Tafeln bestimmen lässt und dass er zugleich vom gegebenen
Logarithmus $-\pi$ nur in der $10^{\mathrm{ten}}$ Dezimalstelle abweichen kann.
Die Zahl, deren Logarithmus diese Differenz ist, musste also nur berechnet
und mit $N$ multipliziert werden, um $A$ zu erhalten,
\[
A=Ne^{\log11+10\log10-\log2281-\log2144-\log972-\pi}=Ne^\delta
\]

\Needspace{34\baselineskip}
\SourceLeaf{436}

Aus der \textsc{Wolframschen} Tafel ergab sich:

[Zur Abkürzung sind $\delta-\log(1+2135.10^{-13})=\delta'$,
\hfill\break{} $\delta'-\log(1+1443.10^{-17})=\delta''$ gesetzt.]

\begingroup\footnotesize\setlength{\arraycolsep}{1pt}\renewcommand{\arraystretch}{1.12}\[
\begin{array}{rcrlllll}
10\log10&&&&&&&\\[5pt]
&=&\GaussIIITranscriptionDigits{23,0258509299}&\GaussIIITranscriptionDigits{4045684017}&\GaussIIITranscriptionDigits{9914546843}&\GaussIIITranscriptionDigits{6420760110}&\GaussIIITranscriptionDigits{14886288}&\\[5pt]
\log11&=&\GaussIIITranscriptionDigits{2,3978952727}&\GaussIIITranscriptionDigits{9837054406}&\GaussIIITranscriptionDigits{1943577965}&\GaussIIITranscriptionDigits{1292998217}&\GaussIIITranscriptionDigits{06853937}&\\[5pt]
&&\GaussIIITranscriptionDigits{25,4237462027}&\GaussIIITranscriptionDigits{3882738424}&\GaussIIITranscriptionDigits{1858124808}&\GaussIIITranscriptionDigits{7713758327}&\GaussIIITranscriptionDigits{21740225}&\\[5pt]
l.2281&=&\GaussIIITranscriptionDigits{7,7323692222}&\GaussIIITranscriptionDigits{8438803081}&\GaussIIITranscriptionDigits{0466064812}&\GaussIIITranscriptionDigits{2168095619}&\GaussIIITranscriptionDigits{53159812}&\\[5pt]
l.2144&=&\GaussIIITranscriptionDigits{7,6704285221}&\GaussIIITranscriptionDigits{9069260675}&\GaussIIITranscriptionDigits{6232603654}&\GaussIIITranscriptionDigits{6055909444}&\GaussIIITranscriptionDigits{33575023}&\\[5pt]
l.972&=&\GaussIIITranscriptionDigits{6,8793558044}&\GaussIIITranscriptionDigits{6043907581}&\GaussIIITranscriptionDigits{0690427528}&\GaussIIITranscriptionDigits{9816593884}&\GaussIIITranscriptionDigits{53057834}&\\[5pt]
\pi&=&\GaussIIITranscriptionDigits{3,1415926535}&\GaussIIITranscriptionDigits{8979323846}&\GaussIIITranscriptionDigits{2643383279}&\GaussIIITranscriptionDigits{5028841971}&\GaussIIITranscriptionDigits{69399375}&\\[5pt]
&&\GaussIIITranscriptionDigits{25,4237462025}&\GaussIIITranscriptionDigits{2531295184}&\GaussIIITranscriptionDigits{0032479275}&\GaussIIITranscriptionDigits{3069440920}&\GaussIIITranscriptionDigits{09192044}&\\[5pt]
\delta&=&\GaussIIITranscriptionDigits{0,.........2}&\GaussIIITranscriptionDigits{1351443240}&\GaussIIITranscriptionDigits{1825645533}&\GaussIIITranscriptionDigits{4644317407}&\GaussIIITranscriptionDigits{12548181}&\\[5pt]
\delta'&=&\GaussIIITranscriptionDigits{0,..........}&\GaussIIITranscriptionDigits{...1443242}&\GaussIIITranscriptionDigits{4616770530}&\GaussIIITranscriptionDigits{2204949495}&\GaussIIITranscriptionDigits{65316893}&\\[5pt]
\delta''&=&\GaussIIITranscriptionDigits{0,..........}&\GaussIIITranscriptionDigits{.......242}&\GaussIIITranscriptionDigits{4616770634}&\GaussIIITranscriptionDigits{3329449495}&\GaussIIITranscriptionDigits{6431533124}&\\[5pt]
\frac12\delta''\delta''&=&\GaussIIITranscriptionDigits{0,..........}&\GaussIIITranscriptionDigits{..........}&\GaussIIITranscriptionDigits{..........}&\GaussIIITranscriptionDigits{.....29393}&\GaussIIITranscriptionDigits{8324222063}&\\[5pt]
e^{\delta''}&=&\GaussIIITranscriptionDigits{1,..........}&\GaussIIITranscriptionDigits{.......242}&\GaussIIITranscriptionDigits{4616770634}&\GaussIIITranscriptionDigits{3329478889}&\GaussIIITranscriptionDigits{4755755187}&\\[5pt]
e^{\delta'}&=&\GaussIIITranscriptionDigits{1,..........}&\GaussIIITranscriptionDigits{...1443242}&\GaussIIITranscriptionDigits{4616770634}&\GaussIIITranscriptionDigits{3679351089}&\GaussIIITranscriptionDigits{4781097632}&\\[5pt]
e^\delta&=&\GaussIIITranscriptionDigits{1,.........2}&\GaussIIITranscriptionDigits{1351443242}&\GaussIIITranscriptionDigits{4619851957}&\GaussIIITranscriptionDigits{0236156393}&\GaussIIITranscriptionDigits{8536512211}&\\[5pt]
11A&=&\GaussIIITranscriptionDigits{0,4753531009}&\GaussIIITranscriptionDigits{0149474751}&\GaussIIITranscriptionDigits{8595108889}&\GaussIIITranscriptionDigits{0081240330}&\GaussIIITranscriptionDigits{0920791693}&\GaussIIITranscriptionDigits{5}\\[5pt]
A&=&\GaussIIITranscriptionDigits{0,0432139182}&\GaussIIITranscriptionDigits{6377224977}&\GaussIIITranscriptionDigits{4417737171}&\GaussIIITranscriptionDigits{7280112757}&\GaussIIITranscriptionDigits{2810981063}&
\end{array}
\]\endgroup

\Needspace{20\baselineskip}
Um die Richtigkeit dieses Ergebnisses durch eine zweite Rechnung
zu bestätigen, multiplizierte man die Zahl $A$ mit $599.10^8$, was
$2588513703,999957..$ ergab. Somit blieb eine sehr einfache Rechnung übrig,
wenn sich die Zahl $2588513704$ in Faktoren kleiner als $10000$ zerlegen ließ.
Ein Versuch ergab $2588513704=8.7.17.2719027$. Es kam
also darauf an, ob $2719027$ eine Primzahl ist. Man fand, dass $-1848$ sicher
ein quadratischer Rest von $2719027$ sein muss, wenn diese Zahl eine Primzahl
ist, und dass sie in diesem Fall \textit{einmal} in der Form $3xx+616yy$ enthalten
sein muss; umgekehrt muss sie durch diese Form entweder gar nicht oder
\textit{mehr als einmal} darstellbar sein, wenn sie zusammengesetzt ist.
Die Ausschlussmethode zeigte jedoch, dass diese Zahl tatsächlich nur einmal
in der Form $3xx+616yy$, nämlich $2719027=3.197^2+616.65^2$,
\SourceLeaf{437}
enthalten ist. Daraus folgte mit Sicherheit, dass $2719027$ eine Primzahl und die
versuchte Methode daher diesmal nicht anwendbar ist.

\Needspace{18\baselineskip}
Nach einigen weiteren vergeblichen Versuchen kam man schließlich auf den
folgenden, der besser gelang. Die Zahl $A$ multipliziert mit $10^{14}$ ergab $4321391826377,25$
und die Zahl $4321391826375$ zerfiel in die Faktoren $125.81.13.32831087$. Dass
die Zahl $32831087$ keine Primzahl ist, folgte daraus, dass sie nicht in der
Form $xx+190yy$ enthalten war. Man wandte daher die erste der in den \textit{Disqui. Arr.}
gelehrten Methoden auf sie an und entdeckte, dass sie das Produkt
aus $373.8819$ ist.

\Needspace{30\baselineskip}
\begin{center}\textit{Zweite Rechnung für $e^{-\pi}=A$}\end{center}

$59.10^9A$ ergab sich sehr genau $=2549621178=54.13.3631939$. Die Zahl
$3631939$ zerfiel in die Faktoren $1091.3329$ (die aus
\[
3631939=40.232^2+19.279^2=40.300^2+19.41^2
\]
gefunden worden waren). Also
\[
\begin{aligned}59A&=702.1091.3329.10^{-9}\\[7pt]&\quad{}e^{9\log10+\log59-\log702-\log1091-\log3329-\pi}=Ne^\delta\end{aligned}
\]
\Needspace{29\baselineskip}
Aus \textsc{Wolframs} Tafel ergab sich [es sind $\delta-\log(1-17157.10^{-14})=\delta'$,
$\frac12\delta'\delta'+\frac16\delta'^3=\varepsilon$ gesetzt]
\begingroup\footnotesize\setlength{\arraycolsep}{1pt}\renewcommand{\arraystretch}{1.12}\[
\begin{array}{rcrlllll}
C9l.10&=&\GaussIIITranscriptionDigits{21,2767341630}&\GaussIIITranscriptionDigits{5358884383}&\GaussIIITranscriptionDigits{8076907840}&\GaussIIITranscriptionDigits{7221315900}&\GaussIIITranscriptionDigits{86602341}&\\[5pt]
C.l.59&=&\GaussIIITranscriptionDigits{5,9224625560}&\GaussIIITranscriptionDigits{9428054938}&\GaussIIITranscriptionDigits{3949626280}&\GaussIIITranscriptionDigits{3023759366}&\GaussIIITranscriptionDigits{53210670}&\\[5pt]
l.702&=&\GaussIIITranscriptionDigits{6,5539334040}&\GaussIIITranscriptionDigits{2581111965}&\GaussIIITranscriptionDigits{6455273791}&\GaussIIITranscriptionDigits{0722868232}&\GaussIIITranscriptionDigits{39752589}&\\[5pt]
l.1091&=&\GaussIIITranscriptionDigits{6,9948499858}&\GaussIIITranscriptionDigits{3307081851}&\GaussIIITranscriptionDigits{1895817110}&\GaussIIITranscriptionDigits{3840806982}&\GaussIIITranscriptionDigits{08318537}&\\[5pt]
l.3329&=&\GaussIIITranscriptionDigits{8,1104272375}&\GaussIIITranscriptionDigits{7502494295}&\GaussIIITranscriptionDigits{2021653586}&\GaussIIITranscriptionDigits{1576658324}&\GaussIIITranscriptionDigits{67001597}&\\[5pt]
\pi&=&\GaussIIITranscriptionDigits{3,1415926535}&\GaussIIITranscriptionDigits{8979323846}&\GaussIIITranscriptionDigits{2643383279}&\GaussIIITranscriptionDigits{5028841971}&\GaussIIITranscriptionDigits{69399375}&\\[5pt]
-\delta&=&\GaussIIITranscriptionDigits{0,.........1}&\GaussIIITranscriptionDigits{7156951280}&\GaussIIITranscriptionDigits{5042661888}&\GaussIIITranscriptionDigits{1414250778}&\GaussIIITranscriptionDigits{24285109}&\\[5pt]
\delta'&=&\GaussIIITranscriptionDigits{0,..........}&\GaussIIITranscriptionDigits{.....48720}&\GaussIIITranscriptionDigits{9675470563}&\GaussIIITranscriptionDigits{5420349120}&\GaussIIITranscriptionDigits{2333831382}&\\[5pt]
\varepsilon&=&\GaussIIITranscriptionDigits{0,..........}&\GaussIIITranscriptionDigits{..........}&\GaussIIITranscriptionDigits{..........}&\GaussIIITranscriptionDigits{1186866339}&\GaussIIITranscriptionDigits{3606594227}&\\[5pt]
e^{\delta'}&=&\GaussIIITranscriptionDigits{1,..........}&\GaussIIITranscriptionDigits{.....48720}&\GaussIIITranscriptionDigits{9675470563}&\GaussIIITranscriptionDigits{6607215459}&\GaussIIITranscriptionDigits{5940425609}&\\[5pt]
1-e^\delta&=&\GaussIIITranscriptionDigits{0,.........1}&\GaussIIITranscriptionDigits{7156951279}&\GaussIIITranscriptionDigits{0324613026}&\GaussIIITranscriptionDigits{9032989386}&\GaussIIITranscriptionDigits{4786573955}&\\[5pt]
59A&=&\GaussIIITranscriptionDigits{2,5496211775}&\GaussIIITranscriptionDigits{6256273669}&\GaussIIITranscriptionDigits{0646493131}&\GaussIIITranscriptionDigits{9526652679}&\GaussIIITranscriptionDigits{5847882752}&\GaussIIITranscriptionDigits{6}\\[5pt]
A&=&\GaussIIITranscriptionDigits{0,0432139182}&\GaussIIITranscriptionDigits{6377224977}&\GaussIIITranscriptionDigits{4417737171}&\GaussIIITranscriptionDigits{7280112757}&\GaussIIITranscriptionDigits{2810981063}&\GaussIIITranscriptionDigits{6}
\end{array}
\]\endgroup

\Needspace{9\baselineskip}
\SourceLeaf{438}

Zugleich lässt sich aus der Übereinstimmung beider Rechnungen schließen,
dass die Logarithmen von $2,3,5,11,13,59,67,1091,2281,3329$ in der
\textsc{Wolframschen} Tafel bis auf die letzte Ziffer richtig sind.

\clearpage
\endgroup
